\providecommand{\Z}{\mathbb{Z}}
\providecommand{\R}{\mathbb{R}}
\providecommand{\C}{\mathbb{C}}
\providecommand{\Lat}{\Lambda}
\providecommand{\cC}{\mathcal{C}}
\providecommand{\cP}{\mathcal{P}}
\providecommand{\cH}{\mathcal{H}}
\providecommand{\cO}{\mathcal{O}}
\providecommand{\fhp}{\mathcal{H}_+^{\mathrm{IV}}}
\providecommand{\Dfive}{\Delta_{5}}
\providecommand{\PhiTenOP}{\Phi_{\!10}^{\mathrm{OP}}}
\providecommand{\Zgeq}{\mathbb{Z}_{\geq 0}}
\providecommand{\BKM}{\mathrm{BKM}}
\providecommand{\CoHA}{\mathrm{CoHA}}
\providecommand{\kch}{\kappa_{\mathrm{ch}}}
\providecommand{\kcat}{\kappa_{\mathrm{cat}}}
\providecommand{\kBKM}{\kappa_{\mathrm{BKM}}}
\providecommand{\kfib}{\kappa_{\mathrm{fiber}}}
\providecommand{\Eone}{E_1}
\providecommand{\Etwo}{E_2}

\section*{The reflection subgroups
\texorpdfstring{$W_I^{(2)}$}{W-I-2} and
\texorpdfstring{$W_{II}^{(2)}$}{W-II-2}}

Let
\[
  \Lat^{2,1}
  = \Lat^{1,1}\oplus[2]
  = \Z f_2 \oplus \Z f_3 \oplus \Z f_{-2},
  \qquad
  (f_2,f_{-2})=-1,\quad (f_3,f_3)=2,
\]
with $f_2$ and $f_{-2}$ isotropic and all other basis pairings zero.
The positive cone
\[
  \cC(\Lat^{2,1})_+
  \subset \{x\in \Lat^{2,1}\otimes\R : (x,x)<0\}
\]
gives the tube chart
\[
  \Omega(\cC(\Lat^{2,1})_+)
  = \Lat^{2,1}\otimes\R+i\cC(\Lat^{2,1})_+
  \subset \fhp .
\]
The inclusion $O(\Lat^{2,1})\hookrightarrow O(\Lat^{3,2})$ extends an
isometry by the identity on $(\Lat^{2,1})^\perp$. The square-two
reflection subgroups, the Maass character of $\Dfive$, and the exact
normalisation of the $N=1$ Borcherds product are the invariants of this
rank-three cusp.

\subsection*{Square-two reflection theorem}

For $\alpha\in\Lat^{2,1}$ with $(\alpha,\alpha)>0$ and
$(\alpha,\alpha)\mid 2(\Lat^{2,1},\alpha)$, the reflection is
\[
  s_\alpha(x)
  =x-\frac{2(x,\alpha)}{(\alpha,\alpha)}\alpha .
\]
For $(\alpha,\alpha)=2$, this is $s_\alpha(x)=x-(x,\alpha)\alpha$.
Set
\[
  \Delta^{(k)}(\Lat^{2,1})
  =\{\delta\in\Lat^{2,1}:(\delta,\delta)=k,\ \delta
  \text{ primitive}\},
\]
and let $W^{(k)}$ be generated by the reflections in
$\Delta^{(k)}(\Lat^{2,1})$.

\paragraph{Theorem.}
Nikulin's rank-three calculation gives
\[
  O(\Lat^{2,1})_+=W^{(2)}(\Lat^{2,1}).
\]

The square-two primitive roots split into two divisor classes:
\[
  \text{Type I}:\quad (\delta,\Lat^{2,1})=\Z,
  \qquad
  \text{Type II}:\quad (\delta,\Lat^{2,1})=2\Z.
\]
The corresponding sublattices are
\[
  \Lat_I^{2,1}
  =\{m f_2+\ell f_3+n f_{-2}:m+\ell+n\equiv0\pmod2\},
\]
\[
  \Lat_{II}^{2,1}
  =\{m f_2+\ell f_3+n f_{-2}:m\equiv n\equiv0\pmod2\}.
\]
Their lattice indices are
$[\Lat^{2,1}:\Lat_I^{2,1}]=2$ and
$[\Lat^{2,1}:\Lat_{II}^{2,1}]=4$. The reflection-subgroup indices are
\[
  [W^{(2)}(\Lat^{2,1}):W^{(2)}(\Lat_I^{2,1})]=2,
  \qquad
  [W^{(2)}(\Lat^{2,1}):W^{(2)}(\Lat_{II}^{2,1})]=6.
\]

\subsection*{Fundamental roots and semidirect products}

The three Coxeter chambers are cut out by square-two roots:
\[
  \cP'=\bigcap_{\delta\in\cP'_{\mathrm{prim}}}\cH^+_\delta.
\]
The primitive walls and chamber automorphism groups are
\begin{align*}
  \cP_{\mathrm{prim}}
    &=\{f_2-f_3,\ f_{-2}-f_2,\ f_3\},
  & A(\cP)&=\langle1\rangle,\\
  \cP_{I,\mathrm{prim}}
    &=\{f_2-f_3,\ f_{-2}-f_2,\ f_2+f_3\},
  & A(\cP_I)&=\langle s_{f_3}\rangle\simeq S_2,\\
  \cP_{II,\mathrm{prim}}
    &=\{\delta_1=2f_2-f_3,\ \delta_2=2f_{-2}-f_3,\ \delta_3=f_3\},
  & A(\cP_{II})&=\langle s_{f_2-f_3},s_{f_{-2}-f_2}\rangle\simeq S_3.
\end{align*}
Direct pairing gives
\[
  ((\delta_i,\delta_j))_{1\leq i,j\leq3}
  =
  \begin{pmatrix}
  2&-2&-2\\
  -2&2&-2\\
  -2&-2&2
  \end{pmatrix}.
\]
Indeed $(\delta_i,\delta_i)=2$,
$(\delta_1,\delta_2)=4(f_2,f_{-2})+(f_3,f_3)=-2$, and
$(\delta_1,\delta_3)=(\delta_2,\delta_3)=-2$.

The chamber decompositions are
\[
  O(\Lat^{2,1})_+
  \simeq W^{(2)}(\Lat^{2,1})
  \simeq W^{(2)}(\Lat_I^{2,1})\rtimes A(\cP_I)
  \simeq W^{(2)}(\Lat_{II}^{2,1})\rtimes A(\cP_{II}).
\]
Under the standard identification $O(\Lat^{2,1})_+\simeq
\mathrm{PGL}_2(\Z)$, the last line is
\[
  \mathrm{PGL}_2(\Z)
  \simeq W^{(2)}(\Lat_{II}^{2,1})\rtimes S_3 .
\]

\subsection*{Maass character of \texorpdfstring{$\Dfive$}{Delta5}}

The Maass multiplier separates the two divisor classes:
\[
  \Dfive(s_\alpha Z)=+\Dfive(Z)
  \quad
  \text{for }
  \alpha\in\{f_{-2}-f_2,\ f_2-f_3,\ f_2+f_3\},
\]
and
\[
  \Dfive(s_{\delta_i}Z)=-\Dfive(Z),
  \qquad i=1,2,3.
\]
\paragraph{Proposition.}
The form $\Dfive$ is invariant under the Type-I generators and
anti-invariant under the Type-II generators. The matrix calculation is
the $\wedge^2$ realisation
\[
  \wedge^2 U=
  \begin{pmatrix}
  0&U\\
  {}^tU^{-1}&0
  \end{pmatrix},
\]
with representatives
\[
  s_{f_{-2}-f_2}
  =-\begin{pmatrix}1&0\\0&1\end{pmatrix},
  \quad
  s_{f_3}
  =-\begin{pmatrix}0&-1\\1&0\end{pmatrix},
  \quad
  s_{f_2-f_3}
  =-\begin{pmatrix}0&-1\\1&-1\end{pmatrix},
\]
followed by Maass' explicit multiplier formula. Thus the character is
carried by the anti-invariant factor in either semidirect decomposition:
\[
  \Dfive(w\,a\,Z)=\det(a)\Dfive(Z)
  \quad
  (w\in W^{(2)}(\Lat_I^{2,1}),\ a\in A(\cP_I)),
\]
\[
  \Dfive(w\,a\,Z)=\det(w)\Dfive(Z)
  \quad
  (w\in W^{(2)}(\Lat_{II}^{2,1}),\ a\in A(\cP_{II})).
\]

\subsection*{Borcherds weight and scalar normalisation}

The reflection computation supplies the chamber/sign component of the
$N=1$ denominator identity. The weight is fixed independently by the
Borcherds weight theorem:
\[
  \kBKM(\Phi_N)=\frac{c_N(0)}{2},
  \qquad N\in\{1,2,3,4,6\}.
\]
For the CHL ladder,
\[
  (c_1(0),c_2(0),c_3(0),c_4(0),c_6(0))
  =(10,8,6,4,2),
\]
hence
\[
  (\kBKM(\Phi_1),\kBKM(\Phi_2),\kBKM(\Phi_3),
  \kBKM(\Phi_4),\kBKM(\Phi_6))
  =(5,4,3,2,1).
\]
At $N=1$, $\Phi_1=\Dfive$, $c_1(0)=10$, and
\[
  \kBKM(\Dfive)=\frac{10}{2}=5.
\]
The primitive Jacobi-form input and the protected K3 elliptic genus obey
\[
  Z^{K3}_{\mathrm{ell}}=2\,\phi_{0,1},
  \qquad
  \operatorname{Borch}(\phi_{0,1})=D_5=64^{-1}\Dfive .
\]
The square produced from the protected index is the
Oberdieck--Pandharipande scalar branch:
\[
  D_5=64^{-1}\Dfive,
  \qquad
  \PhiTenOP
  =D_5^2
  =4096^{-1}\Dfive^2.
\]
The scalar $4096=64^2$ supplies the OP square normalisation. Hall
orientation characters require source orientation-line monodromy in the
critical CoHA geometry.

\subsection*{Separation of invariants and output levels}

For $K3\times E$ the compact categorical and compact chiral invariants
share the total-space value, while the Heisenberg, Borcherds, and
fibre constructions give separate numbers:
\[
  \kcat(K3\times E)
  =\kch(K3\times E)
  =\chi(\cO_{K3})\chi(\cO_E)
  =2\cdot0=0,
\]
\[
  \kappa_{\mathrm{ch}}^{\mathrm{Heis}}(K3\times E)=3,
  \qquad
  \kBKM(\mathfrak g_{\Dfive})=5,
  \qquad
  \kfib(K3)=24.
\]
Thus the displayed spectrum of values is $\{0,3,5,24\}$, with the
zero shared by two compact total-space invariants.

For $d\geq3$, the native output of $\Phi_d$ is $\Eone$-chiral. A braided
$\Etwo$ structure appears, where constructed, on
$\mathcal Z(\mathrm{Rep}^{\Eone}(A))$; the algebra $A$ itself remains
in the native $\Eone$ layer.
The toric comparison obeys the same separation:
\[
  \CoHA(\C^3)=Y^+(\widehat{\mathfrak{gl}}_1),
\]
the positive half. The full $\mathcal W_{1+\infty}$ object is reached
only after Drinfeld double or centre passage and vacuum-module
evaluation. The six routes to $G(K3\times E)$ are six distinct
constructions of comparison objects for the same geometry.

The Hall--Drinfeld double belongs to the Borcherds denominator branch.
The K3 self-mirror
Yangian branch belongs to the separate Mukai-lattice presentation.

\subsection*{Weyl vector}

The three Type-II simple roots determine
\[
  \rho=\frac12(\delta_1+\delta_2+\delta_3)
  =\frac12\bigl((2f_2-f_3)+(2f_{-2}-f_3)+f_3\bigr)
  =f_2-\frac12 f_3+f_{-2}.
\]
This $\rho$ satisfies $(\rho,\delta_i)=-1$ for
$i=1,2,3$, hence it is the Weyl vector of the rank-three real-root
subsystem. The equality
\[
  \frac{1}{64}\Dfive(2Z)
  =
  e^{-2\pi i(\rho,z)}
  \prod_{\alpha\in\Delta_+}
  (1-e^{-2\pi i(\alpha,z)})^{\mathrm{mult}\,\alpha}
\]
is the Borcherds--Gritsenko--Nikulin denominator identity in the
$0$-cusp tube chart. The Weyl-vector calculation checks the real-root
Cartan and exponential prefactor; the numerical value
$\kBKM(\Dfive)=5$ comes from $c_1(0)/2$. Quadratic expressions in
$\rho$ control the chamber prefactor.

\subsection*{Proof obligations}

\begin{enumerate}
\item The reflection/sign lemma is a target-side automorphic statement
after citing Gritsenko--Nikulin and Maass. Source-side Hall orientation
characters require an independent orientation-line calculation.
\item Any Hall or CoHA theorem using the Type-II sign must compute the
orientation-line monodromy on the source moduli problem: local wall
charts, reduced Ext normal forms, invariant unit character, and
Pfaffian rank.
\item Any compact $K3\times E$ Hall--Borcherds identification must
construct the oriented critical CoHA positive half, its nondegenerate
Hopf or stable-envelope pairing, and the Drinfeld double. The scalar
OP identity $Z^{K3\times E}_{\mathrm{OP}}=-4096\,\Dfive^{-2}$ supplies
enumerative evidence for the denominator square.
\item Any use of the $S_3$ semidirect decomposition and the
Type-I/Type-II character formula requires the separate normalisation line
$D_5=64^{-1}\Dfive$ and
$\PhiTenOP=4096^{-1}\Dfive^2$, so that the Borcherds
weight, the monic product, and the OP scalar branch remain separate.
\end{enumerate}
